\DocumentMetadata{
  pdfversion=2.0,
  pdfstandard=ua-2,
  lang=en-US,
  tagging=on,
  tagging-setup={math/setup=mathml-SE,tabsorder=structure}
}
\documentclass[11pt,letterpaper]{article}
\usepackage[margin=0.68in,includefoot]{geometry}
\usepackage{newcomputermodern}
\usepackage{amsmath,amscd,mathtools}
\providecommand{\square}{\mdlgwhtsquare}
\usepackage[hidelinks]{hyperref}
\hypersetup{
  pdftitle={Numerical Analysis PhD Exam, May 2019},
  pdfauthor={University of Florida Department of Mathematics},
  pdfsubject={Graduate mathematics examination},
  pdfdisplaydoctitle=true,
  bookmarksopen=true
}
\setcounter{secnumdepth}{0}
\setlength{\parindent}{0pt}
\setlength{\parskip}{0.28em}
\setlength{\emergencystretch}{3em}
\setlength{\leftmargini}{2.15em}
\setlength{\leftmarginii}{2.65em}
\setlength{\leftmarginiii}{3em}
\renewcommand{\labelenumi}{\textbf{\theenumi.}}
\pagestyle{empty}
\raggedbottom
\allowdisplaybreaks
\newenvironment{examproblems}[1]{%
  \begin{enumerate}%
  \setcounter{enumi}{#1}%
  \setlength{\topsep}{0.35em}%
  \setlength{\partopsep}{0pt}%
  \setlength{\itemsep}{0.48em}%
  \setlength{\parsep}{0.18em}%
}{\end{enumerate}}
\newenvironment{examsubparts}{%
  \begin{enumerate}%
  \setlength{\topsep}{0.18em}%
  \setlength{\partopsep}{0pt}%
  \setlength{\itemsep}{0.16em}%
  \setlength{\parsep}{0.1em}%
}{\end{enumerate}}
\newenvironment{examinstructions}{%
  \par\begingroup\itshape
}{%
  \par\endgroup\vspace{0.18em}
}
\makeatletter
\renewcommand\section{\@startsection{section}{1}{0pt}%
  {-1.25ex plus -.25ex minus -.1ex}%
  {0.55ex plus .1ex}%
  {\normalfont\large\bfseries}}
\renewcommand{\@maketitle}{%
  \newpage\null\vskip -1.2em%
  {\footnotesize\bfseries
    \href{https://www.ufl.edu/}{University of Florida}, \href{https://math.ufl.edu/}{Department of Mathematics}\par}%
  \vskip 0.18em%
  \tagstructbegin{tag=Title}%
    {\LARGE\bfseries\href{https://gma.math.ufl.edu/exams/numerical-analysis-phd/}{Numerical Analysis PhD Exam}, May 2019\par}%
  \tagstructend%
  \vskip 0.35em%
  \tagmcbegin{artifact}%
    \hrule height 1pt%
  \tagmcend%
  \vskip 0.55em%
}
\makeatother
\title{Numerical Analysis PhD Exam, May 2019}
\author{}
\date{}
\begin{document}
\maketitle
\thispagestyle{empty}
\section{Numerical Linear Algebra}
\begin{examinstructions}
Do 4 (four) problems.
\end{examinstructions}
\begin{examproblems}{0}
\item
This problem has three parts.
\begin{examsubparts}
\item[\textnormal{(a)}]
Show the matrix norm equality for \(A\in\mathbb{C}^{m\times n}\) 
\[
\lVert A\rVert_\infty=\max_{1\le i\le m}\sum_{j=1}^n |a_{ij}|.
\]
\item[\textnormal{(b)}]
Explain why the matrix 1-norm, 2-norm and~\(\infty\)-norm are the most commonly used of the matrix~\(p\)-norms in scientific computing.
\item[\textnormal{(c)}]
Show \(\rho(A)\le \lVert A\rVert\), where \(\lVert A\rVert\) is any subordinate (induced) matrix norm and \(\rho(A)\) is the spectral radius of~\(A\).
\end{examsubparts}
\item
Let \(A\in\mathbb{C}^{m\times n}\) with \(\operatorname{rank}(A)=n<m\). Let \(A=QR\) be the \(QR\) decomposition of~\(A\), and \(A=Q_1R_1\) be the economy \(QR\) decomposition.
\begin{examsubparts}
\item[\textnormal{(a)}]
Show \(Q_1Q_1^*\) is an orthogonal projector onto \(\operatorname{Col}(A)\).
\item[\textnormal{(b)}]
Let \(b\in\mathbb{C}^m\). Write down an expression for the least-squares solution to \(Ax=b\) as the solution to an \(n\times n\) system in terms of \(Q_1\), (and/or \(Q_1^*\)), \(R_1\), \(x\) and~\(b\).
\end{examsubparts}
\item
Let \(A=U\Sigma V^*\) be the singular value decomposition of \(A\in\mathbb{C}^{m\times n}\) with \(\operatorname{rank}(A)=p\) and \(p\le n\le m\).
\begin{examsubparts}
\item[\textnormal{(a)}]
Show \(\operatorname{Col}(A^*)=\operatorname{Span}\{v_1,v_2,\ldots,v_p\}\), where \(v_1,\ldots,v_p\) are the first~\(p\) columns of~\(V\).
\item[\textnormal{(b)}]
Show \(\operatorname{Null}(A)=\operatorname{Span}\{v_{p+1},v_{p+2},\ldots,v_n\}\).
\item[\textnormal{(c)}]
Suppose the right singular vectors \(v_1,\ldots,v_p\) have been computed. Describe how to compute the left singular vectors \(u_1,\ldots,u_p\) (without solving a spectral problem).
\end{examsubparts}
\item
Let \(\lVert\cdot\rVert\) be a subordinate (induced) matrix norm. If~\(A\) is \(n\times n\) invertible and~\(E\) is \(n\times n\) with \(\lVert A^{-1}\rVert\lVert E\rVert<1\), then show
\begin{examsubparts}
\item[\textnormal{(a)}]
\(A+E\) is nonsingular.
\item[\textnormal{(b)}]
\[
\lVert(A+E)^{-1}\rVert\le \frac{\lVert A^{-1}\rVert}{1-\lVert A^{-1}\rVert\lVert E\rVert}.
\]
\end{examsubparts}
\item
Consider the matrix~\(A\) given by 
\[
A=\begin{pmatrix}
1&-1&2&0\\
-1&4&-1&1\\
2&-1&6&-2\\
0&1&-2&4
\end{pmatrix}.
\]
 Suppose the eigenvalues of~\(A\) are all distinct (they are) and satisfy \(\lambda_1>\lambda_2>\lambda_3>\lambda_4\).
\begin{examsubparts}
\item[\textnormal{(a)}]
Show that~\(A\) is positive definite.
\item[\textnormal{(b)}]
Describe an algorithm that could be used to converge to \(\lambda_4\).
\item[\textnormal{(c)}]
Describe an algorithm that could be used to converge to \(\lambda_2\).
\end{examsubparts}
\end{examproblems}
\section{Numerical Analysis}
\begin{examinstructions}
Do 4 (four) problems.
\end{examinstructions}
\begin{examproblems}{0}
\item
Consider the fixed point problem \(x=f(x)\), where \(f(x)=e^{-(2+x)}\).
\begin{examsubparts}
\item[\textnormal{(a)}]
Find the largest open interval in~\(\mathbb{R}\) where \(f(x)\) is a contraction.
\item[\textnormal{(b)}]
Assuming all computations are done in exact arithmetic, find the largest open interval in~\(\mathbb{R}\) where the fixed-point iteration \(x_{k+1}=f(x_k)\) is assured to converge.
\item[\textnormal{(c)}]
Write a Newton iteration for finding the fixed-point.
\end{examsubparts}
\item
Let \(x_1,x_2,\ldots,x_{n+1}\) be \(n+1\) distinct numbers. Let \(l_j(x)\) be the associated Lagrange basis polynomials, \(j=1,\ldots,n+1\).
\begin{examsubparts}
\item[\textnormal{(a)}]
State the definition of \(l_j(x)\) and show that \(\{l_j(x)\}_{j=1}^{n+1}\) form a basis for \(P_n\), the space of polynomials of degree at most~\(n\).
\item[\textnormal{(b)}]
Show that 
\[
\sum_{j=1}^{n+1}(x-x_j)^k l_j(x)=0,\qquad \text{for all }k=1,\ldots,n.
\]
\end{examsubparts}
\item
Consider the interval \([a,b]\) with the partition \(a=x_1<x_2<\cdots<x_n<x_{n+1}=b\). Suppose \(s(x)\) is the natural cubic spline that interpolates the data \(\{(x_i,y_i)\}_{i=1}^{n+1}\), and that \(g\in C^2[a,b]\) interpolates the same data. Show that 
\[
\int_a^b \bigl(s''(x)\bigr)^2\,dx\le \int_a^b \bigl(g''(x)\bigr)^2\,dx.
\]
\item
This problem has two parts.
\begin{examsubparts}
\item[\textnormal{(a)}]
Consider the inner product on \(C(0,2)\) given by \((f,g)=\int_0^2 f(t)g(t)\,dt\). Find three orthonormal polynomials \(\phi_0,\phi_1,\phi_2\) on \((0,2)\) with respect to the given inner product such that the degree of \(\phi_n\) is equal to~\(n\), \(n=0,1,2\).
\item[\textnormal{(b)}]
Find the nodes \(t_1\) and \(t_2\) and weights \(w_1\) and \(w_2\) which yield the weighted Gaussian Quadrature formula 
\[
\int_0^2 f(t)\,dt\approx w_1f(t_1)+w_2f(t_2)
\]
 with degree of exactness \(m=3\). You should find the nodes exactly, and may leave the weights \(w_1,w_2\) in integral form.
\end{examsubparts}
\item
Let \(f\in C^\infty(a-H,a+H)\), and let \(h<H\). Let \(x_0=a-h\), \(x_1=a\) and \(x_2=a+h\).
\begin{examsubparts}
\item[\textnormal{(a)}]
Find the finite difference approximation to \(f'(a)\) based on the interpolant \(p_2\) which satisfies \(p_2(x_0)=f(x_0)\), \(p_2(x_1)=f(x_1)\) and \(p_2(x_2)=f(x_2)\).
\item[\textnormal{(b)}]
Let \(\psi_0(h)=\psi(h)\) be the difference approximation to \(f'(a)\) found in part (a). Assume (in exact arithmetic) \(\psi(h)\to\psi(0)=f'(a)\) as \(h\to0\), and that \(\psi(h)\) has the asymptotic expansion 
\[
\psi(h)=\psi(0)+a_2h^2+a_4h^4+a_6h^6+\cdots.
\]
 Find the general Richardson extrapolation formula for \(\psi_k(h)\) based on \(\psi_{k-1}(h)\) and \(\psi_{k-1}(h/2)\).
\end{examsubparts}
\end{examproblems}
\end{document}
